\section{Results from the companion papers}\label{app:companions}
This is the complete companion-import list for the present manuscript. The versioned references are Paper~A~\cite{PaperA}, v2.0.1; Paper~B~\cite{PaperB}, v1.0.0; and Paper~C~\cite{PaperC}, v1.0.0. Numbers and source labels below refer to these versions, whose DOIs are in the bibliography.

\begin{enumerate}
\item \textbf{A, Lemmas 2.5--2.6, \texttt{prelim.tex}, labels \texttt{lem:window}, \texttt{lem:ratspec}}. A supported Weyl vector has residual at most $8/M$; rational phase unions approach the irrational spectrum with a square-root frequency bound.

\emph{Use here:} Equation~\eqref{eq:continuity}; the surplus conversion; the derivative domain; Corollary~\ref{cor:derivatives}.

\item \textbf{A, Proposition 3.3(1), \texttt{polynomial.tex}, label \texttt{prop:poly}}. $P(H)=V+S^q+S^{-q}+R_q$, with range-$q$ self-adjoint remainder bounded by $C_hq^{5/2}e^{hq}|\delta|$.

\emph{Use here:} Equations~\eqref{eq:polyidentity}--\eqref{eq:regroup}; interpretation of the coarse norm error.

\item \textbf{A, Proposition 3.3(2) and proof Steps 1--3}. Near-spectrum derivative bounds $|P^{(r)}|\le C_{r,h}q^re^{hq}$; locally modified rational comparison and a fixed smooth cutoff.

\emph{Use here:} Top-band length and width bounds; inverse curvature; $m_0$; Proposition~\ref{prop:even}; Proposition~\ref{prop:remainder}.

\item \textbf{A, Theorem A}. Typical moment and fastest upper Abel exponents $\max(0,1-L/\beta)$; all-phase statement for $\beta=\infty$. A gives the finite-positive-$\beta$ law used here; its $\beta=\infty$ transport and packing endpoint is due to Jitomirskaya--Zhang~\cite[Corollary~3]{JZ2022}, and its $\beta=0$ fastest endpoint to Damanik--Tcheremchantsev~\cite{DT2007}.

\emph{Use here:} Section~\ref{sec:diagram}, including the typical-transport entry and the all-phase infinite-$\beta$ endpoint.

\item \textbf{A, Theorem C}. Both packing dimensions of every nonzero spectral measure have the same value under its a.e.-phase hypotheses, and at every phase for $\beta=\infty$ (the Jitomirskaya--Zhang endpoint~\cite[Corollary~3]{JZ2022}).

\emph{Use here:} The spectral-measure packing statement in Section~\ref{sec:diagram}.

\item \textbf{B, Lemmas 2.4--2.5, \texttt{prelim.tex}, labels \texttt{lem:growth}, \texttt{lem:ratprod}}. Uniform real-energy transfer growth and its rational-product expansion.

\emph{Use here:} Appendix~\ref{app:complex} proves the needed complex extension. The extension is proved here, rather than quoted as the original statement.

\item \textbf{B, Theorem A}. The typical fastest-exponent law for supercritical type I operators.

\emph{Use here:} Background attribution in Section~\ref{sec:diagram}; the AMO table itself uses A.

\item \textbf{C, Corollary 7.2, \texttt{sharp.tex}, label \texttt{cor:closedgap}}. For even $q$, $P$ is even and $P(0)=(-1)^{q/2}(2\lambda^q+2)$.

\emph{Use here:} Proposition~\ref{prop:even}; the central-cone discussion in Appendix~\ref{app:remainder}.

\item \textbf{C, Theorem 7.5, \texttt{sharp.tex}, label \texttt{thm:quarter}}. At even $q_n$, a polynomially small fastest tail at radius $c_\eta Te^{-hq_n/2}$ for $b_n\ge h/2+(h-L)$ and $e^{(h/2+(h-L))q_n}\le T\le|\delta_n|^{-1}$; nearby turning-phase clause.

\emph{Use here:} The parity-qualified phase-$1/4$ lower bound and transport geometry in Section~\ref{sec:diagram}.

\item \textbf{C, Corollary 5.6, \texttt{upperlaw.tex}, label \texttt{cor:typeI}}. For AMO, $\alpha_u^{\rm unif}\le\beta/(\beta+L)$ below $\beta=L$ and $\le1-L/(2\beta)$ above, including $\beta=L$.

\emph{Use here:} Every upper transport entry in Table~\ref{tab:phase}; the endpoint $\beta=0$.

\item \textbf{C, Corollary 7.6, \texttt{sharp.tex}, label \texttt{cor:exactunif}}. If $\beta_e=\beta\in[L,\infty)$, the uniform, worst-phase and phase-$1/4$ exponents equal $1-L/(2\beta)$.

\emph{Use here:} Exact entries at $x\ge2$ with the stated parity hypothesis.

\item \textbf{C, Proposition 6.4, \texttt{amoproof.tex}, label \texttt{prop:Camo}}. Uniform confinement radius bounded by a constant times $\sqrt{\sigma_n}/|\delta_n|$ plus its specified collars.

\emph{Use here:} The geometric comparison in Section~\ref{sec:diagram}; not a new confinement estimate.
\end{enumerate}

In A v2.0.1, Lemma~2.5 is the window lemma and Lemma~2.6 the rational-spectrum lemma; in C v1.0.0, Corollary~7.2 is the closure corollary and 7.3 its priority remark. The source labels guard against numbering changes between versions.

For clarity, the external analytic inputs are separate from this table. Chambers' formula is due to Chambers~\cite{Chambers1965}; the periodic polynomial identity also follows from the standard magic formula~\cite[Theorem~3.1]{DKS2010}. On the spectrum, the Lyapunov exponent equals $\log\lambda$, by Bourgain--Jitomirskaya~\cite[Corollary~2]{BourgainJitomirskaya2002}, after matching the coupling convention. The IDS facts are standard AMO background, as used in A/B and reviewed in~\cite{AYZdry}. The rational core measure is Avron--van Mouche--Simon~\cite{AvMS}; the all-irrational noncritical measure is Jitomirskaya--Krasovsky~\cite{JKmeasure}. The one-well quantization input is explicitly scoped in Appendix~\ref{app:mathieu}. The Aubry--Perron fact needed at the reference fibre is proved in Lemma~\ref{lem:Wannier}. No unproved general Harper perturbation lemma, all-band reduction, or dynamical localization assumption is imported.

For the new gap theorem the companion inputs are A, Proposition~3.3(2), and B, Lemmas~2.4--2.5, as extended in Appendices~\ref{app:complex} and~\ref{app:below-family}. The one-band range extension is proved in Lemma~\ref{lem:belowrange}; the new inverse, locality, self-energy and model estimates are proved in Appendices~\ref{app:below-resolvent}--\ref{app:below-model}. The additional lower bound uses only Liu--Shi, Theorem~1.3, and the duality proof in Section~\ref{sec:below-literature}. Paper~C supplies the separately labelled transport picture and even-denominator surplus variant. No transport theorem is imported into the new fixed-$D$ proof.

\label{end:companions}
